\section*{Sources and evidence}
\addcontentsline{toc}{section}{Sources and evidence}
Primary access checks on 6 October 2026 cover Boneh's author publication abstract,
the BIPM mole definition, and Landauer's scanned primary reprint, selected
Sections 3--4 and the start of Section 5 (reprint pages 264--266).
No molecular resource measurement or thermodynamic experiment was reproduced.
The occupancy, completion, transport and detection calculations are explicit
teaching models; assigned parameters do not become measured DNA rates.
Independent scientific review remains open.

\begin{thebibliography}{9}
\bibitem{dnad24:boneh} D. Boneh, C. Dunworth, R. J. Lipton and J. Sgall.
\emph{On the computational power of DNA}.
Discrete Applied Mathematics 71 (1996), 79--94.
Author abstract and publication record:
\url{https://crypto.stanford.edu/~dabo/pubs/abstracts/biocircuit.html}.
\bibitem{dnad24:si} Bureau International des Poids et Mesures.
\emph{SI base unit: mole (mol)}. Official SI definition, accessed 6 October 2026.
\url{https://www.bipm.org/en/si-base-units/mole}.
\bibitem{dnad24:landauer} R. Landauer.
\emph{Irreversibility and Heat Generation in the Computing Process}.
IBM Journal of Research and Development 5 (1961), 183--191.
\url{https://doi.org/10.1147/rd.53.0183}.
Primary reprint, IBM Journal of Research and Development 44 (2000), 261--269:
\url{https://www.cpt.univ-mrs.fr/~verga/pdfs/Landauer-1961uq.pdf}.
\end{thebibliography}
